\hypertarget{appendix-2e245a6fa8d3}{}
\section*{Short-circuit evaluation}
\section{Short-circuit evaluation}\label{short-circuit-evaluation}

\textbf{Curriculum depth:} Introduced · \textbf{First seen:} C001 · \textbf{Latest development:} C001

\subsection{Definition and mechanism}\label{definition-and-mechanism}

\texttt{\&\&} may skip its right operand after false; \texttt{\textbar{}\textbar{}} may skip it after true, so evaluation safety is part of semantics.

Built-in logical AND skips its right operand after false; logical OR skips it after true. \texttt{ptr\ \&\&\ ptr-\textgreater{}ready()} guards the dereference. A bitwise \texttt{\&} evaluates both operands. Operator spelling therefore determines both result meaning and whether a dangerous expression is evaluated.

\subsection{Prerequisites}\label{prerequisites}

\begin{itemize}
\tightlist
\item
  \hyperref[appendix-fe6b457d68d8]{Logical versus bitwise operators} --- Built-in \texttt{\&\&}, \texttt{\textbar{}\textbar{}}, and \texttt{!} convert operands to truth conditions.
\end{itemize}

\subsection{Used by}\label{used-by}

\begin{itemize}
\tightlist
\item
  \hyperref[appendix-652d81568de9]{Checked binary parsing} --- uses this mechanism to define its representation, error or validation contract.
\end{itemize}

Connect later chunks here when they use this concept.

\subsection{Recall and next depth}\label{recall-and-next-depth}

Explain short-circuit evaluation without the label, work the example above, then state the assumption whose removal would change the result.

\textbf{Pending:} sequencing, side effects, and optimizer control flow.

\subsection{Source and navigation}\label{source-and-navigation}

\href{Continuous\%20Notes/Computing\%20Foundations\%20-\%20From\%20Bits\%20to\%20C++\%20and\%20Systems\%20APIs\#From\%20a\%20signal\%20to\%20an\%20interpretation}{C001 teaching chapter} · \hyperref[appendix-fd090d982a15]{Part definitions and mechanisms} · \href{}{Dependency Map (02-Foundations-II, appendix-4ad47bdd269a)}
